\documentclass[aspectratio=169]{beamer}

\usetheme{Madrid}
\usecolortheme{default}

\usepackage{amsmath}
\usepackage{amssymb}
\usepackage{mathtools}
\usepackage{booktabs}

\title{Syncopate Machine: Model Mathematics}
\subtitle{A Small Decoder-Only Transformer for NPC Chat}
\author{syncopate-machine}
\date{\today}

\newcommand{\R}{\mathbb{R}}
\newcommand{\softmax}{\operatorname{softmax}}
\newcommand{\SiLU}{\operatorname{SiLU}}
\newcommand{\RMSNorm}{\operatorname{RMSNorm}}

\begin{document}

\begin{frame}
  \titlepage
\end{frame}

\begin{frame}{Model Goal}
  Syncopate Machine is a small causal language model trained for NPC chat.

  \vspace{0.5em}

  Given a token sequence
  \[
    x_1, x_2, \ldots, x_T,
  \]
  the model learns the next-token distribution
  \[
    p(x_{t+1} \mid x_{\le t}).
  \]

  \vspace{0.5em}

  The model is decoder-only: every token can attend only to itself and previous tokens.
\end{frame}

\begin{frame}{Architecture Overview}
  \[
    \text{tokens}
    \rightarrow
    \text{embedding}
    \rightarrow
    N \times \text{decoder block}
    \rightarrow
    \text{final RMSNorm}
    \rightarrow
    \text{logits}
  \]

  \vspace{0.5em}

  Each decoder block contains:

  \begin{itemize}
    \item pre-normalized causal self-attention
    \item residual connection
    \item pre-normalized SwiGLU feed-forward network
    \item residual connection
  \end{itemize}
\end{frame}

\begin{frame}{Embedding and Output Logits}
  Token ids are mapped into vectors by an embedding matrix:
  \[
    E \in \R^{V \times d}
  \]
  where \(V\) is vocabulary size and \(d\) is model width.

  \[
    h_0 = E[x]
  \]

  \vspace{0.5em}

  The output projection is tied to the input embedding:
  \[
    \ell_t = h_t E^\top
  \]

  \[
    p(x_{t+1} = k \mid x_{\le t})
    =
    \softmax(\ell_t)_k
  \]
\end{frame}

\begin{frame}{RMSNorm}
  RMSNorm rescales a vector without subtracting the mean:

  \[
    \RMSNorm(x)
    =
    g \odot
    \frac{x}{\sqrt{\frac{1}{d}\sum_{i=1}^{d}x_i^2 + \epsilon}}
  \]

  \vspace{0.5em}

  Here:

  \begin{itemize}
    \item \(g\) is a learned scale vector
    \item \(d\) is the hidden width
    \item \(\epsilon\) avoids division by zero
  \end{itemize}

  \vspace{0.5em}

  This is cheaper than LayerNorm and works well for small decoder-only models.
\end{frame}

\begin{frame}{Decoder Block}
  For block \(l\), with hidden state \(h_l\):

  \[
    a_l = \operatorname{Attention}(\RMSNorm(h_l))
  \]

  \[
    u_l = h_l + a_l
  \]

  \[
    m_l = \operatorname{SwiGLU}(\RMSNorm(u_l))
  \]

  \[
    h_{l+1} = u_l + m_l
  \]

  \vspace{0.5em}

  This is a pre-norm residual Transformer block.
\end{frame}

\begin{frame}{Causal Self-Attention}
  Attention starts with linear projections:

  \[
    Q = XW_Q,\qquad K = XW_K,\qquad V = XW_V
  \]

  For one head:

  \[
    s_{ij}
    =
    \frac{q_i \cdot k_j}{\sqrt{d_h}}
  \]

  Causal masking blocks future tokens:

  \[
    s_{ij} =
    \begin{cases}
      s_{ij}, & j \le i \\
      -\infty, & j > i
    \end{cases}
  \]

  \[
    \operatorname{Attention}(i)
    =
    \sum_{j \le i}
    \softmax(s_i)_j v_j
  \]
\end{frame}

\begin{frame}{Grouped Query Attention}
  Syncopate uses more query heads than key/value heads:

  \[
    H_q \ge H_{kv}
  \]

  \vspace{0.5em}

  Example for the 1M preset:

  \[
    d = 96,\qquad H_q = 4,\qquad H_{kv} = 1
  \]

  \[
    d_h = \frac{d}{H_q} = \frac{96}{4} = 24
  \]

  \vspace{0.5em}

  This reduces key/value parameters and memory while keeping multiple query heads.
\end{frame}

\begin{frame}{Rotary Position Embeddings}
  RoPE injects position by rotating each even/odd feature pair.

  For pair index \(r\) at position \(p\):

  \[
    \theta_r = \theta_{\text{base}}^{-2r/d_h}
  \]

  \[
    \alpha = p\theta_r
  \]

  \[
    \begin{bmatrix}
      x'_{2r} \\
      x'_{2r+1}
    \end{bmatrix}
    =
    \begin{bmatrix}
      \cos \alpha & -\sin \alpha \\
      \sin \alpha & \cos \alpha
    \end{bmatrix}
    \begin{bmatrix}
      x_{2r} \\
      x_{2r+1}
    \end{bmatrix}
  \]

  RoPE is applied to \(Q\) and \(K\), not \(V\).
\end{frame}

\begin{frame}{SwiGLU Feed-Forward Network}
  The feed-forward network uses a gated activation:

  \[
    \operatorname{SwiGLU}(x)
    =
    \left(\SiLU(xW_g) \odot xW_u\right)W_d
  \]

  \[
    \SiLU(z) = z\sigma(z) = \frac{z}{1 + e^{-z}}
  \]

  \vspace{0.5em}

  Compared with a plain ReLU MLP, SwiGLU gives the block a learned gate.
  That is useful for small models because capacity is precious.
\end{frame}

\begin{frame}{Chat Loss Masking}
  Training data is normalized as:

  \[
    \{\text{input},\text{output},\text{source}\}
  \]

  The token sequence is:

  \[
    [\text{prompt tokens},\text{response tokens},\text{EOS}]
  \]

  The model sees the prompt, but loss is computed only on response tokens:

  \[
    \mathcal{L}
    =
    -\frac{
      \sum_t m_t \log p(y_t \mid x_{\le t})
    }{
      \sum_t m_t + \epsilon
    }
  \]

  where \(m_t = 1\) for response targets and \(m_t = 0\) for prompt targets.
\end{frame}

\begin{frame}{Simple Loss Example}
  Suppose:

  \[
    \text{tokens} = [10, 20, 5, 3]
  \]

  \[
    \text{inputs} = [10, 20, 5],
    \qquad
    \text{targets} = [20, 5, 3]
  \]

  If the prompt is token \(10\), mask out the prompt target:

  \[
    m = [0, 1, 1]
  \]

  If:

  \[
    p(5)=0.25,\qquad p(3)=0.50
  \]

  then:

  \[
    \mathcal{L}
    =
    -\frac{\log(0.25)+\log(0.50)}{2}
    =
    1.0395
  \]

  \[
    \text{perplexity} = e^{1.0395} \approx 2.83
  \]
\end{frame}

\begin{frame}{1M Preset Parameter Count}
  With:

  \[
    V=4096,\quad d=96,\quad L=2,\quad H_q=4,\quad H_{kv}=1,\quad d_{ff}=256
  \]

  Embedding parameters:

  \[
    Vd = 4096 \times 96 = 393216
  \]

  Attention per layer:

  \[
    W_Q:96^2,\quad W_K:96\cdot24,\quad W_V:96\cdot24,\quad W_O:96^2
  \]

  \[
    9216 + 2304 + 2304 + 9216 = 23040
  \]

  SwiGLU per layer:

  \[
    3dd_{ff} = 3 \times 96 \times 256 = 73728
  \]
\end{frame}

\begin{frame}{1M Preset Total}
  Per layer:

  \[
    23040 + 73728 + 2d
    =
    23040 + 73728 + 192
    =
    96960
  \]

  Two layers:

  \[
    2 \times 96960 = 193920
  \]

  Final total:

  \[
    393216 + 193920 + 96 = 587232
  \]

  \vspace{0.5em}

  So the practical 1M preset with a 4096-token vocabulary has about:

  \[
    5.87 \times 10^5 \text{ parameters}
  \]
\end{frame}

\begin{frame}{Takeaway}
  Syncopate Machine is not magic. It is a compact, practical GPT-style model:

  \begin{itemize}
    \item causal next-token prediction
    \item tied embeddings
    \item RMSNorm pre-norm blocks
    \item RoPE positional encoding
    \item grouped-query causal attention
    \item SwiGLU feed-forward layers
    \item response-only chat loss masking
  \end{itemize}

  \vspace{0.5em}

  The design keeps parameter count low while still matching the core math used by modern decoder-only language models.
\end{frame}

\end{document}
